\begin{afterword}
\summaryhead

The post takes its material from Robert Cialdini's \textsc{Influence}. Scarcity, in social psychology, is things becoming more desirable as they seem less obtainable. Two-year-old boys went for a toy behind a tall barrier. After a ban in Miami, residents rated phosphate detergents more highly than people in Tampa did. Information behaves the same way: students grew more opposed to coed dorms when a speech against them was banned, jurors told to disregard a driver's insurance awarded more, and supermarket buyers told of a beef shortage ordered twice as much, or six times as much when told the news was exclusive.

The post explains this by psychological reactance: people resist restrictions on their freedom and rush to seize options that are disappearing. That may have suited hunter-gatherers but is costly with money, as an appliance salesman's trick of saying the last one was just sold shows. Cialdini's warning sign is wanting to possess a thing rather than use it. The deeper problem is that whatever is obtained stops being unattainable, so a person who cannot enjoy what is available will always be frustrated.

\responsehead

The post reports its examples accurately from its source. The toddler, detergent, jury and beef figures all match Cialdini's account. The problems are in what the post leaves out of that source, and in how firmly it presents the studies.

The post presents the scarcity response only as a fault: perhaps an old adaptation, ``rather costly'' in a money economy, a ``malfunction.'' Its source is more even-handed. Cialdini gives two causes of the effect, and the first is a shortcut that usually works: things that are hard to get are typically better, so availability is a fair guide to quality, and ``by following it, we are usually and efficiently right.'' The post names only the second cause, reactance, the urge to resist a restriction on one's freedom. Cialdini also separates honest scarcity from the faked kind. The beef study, the post's most striking example, is the honest kind; by the post's own parenthesis, and by Cialdini's note, the shortage was real. Buyers who ordered more on correct news of a coming shortage are not shown to have made a mistake. The appliance trick, where the scarcity is faked, is the case that fits the post's moral.

The post also changes the remedy. Cialdini does not treat wanting to possess something as a sign of error. The book says that if you want a rare thing for the benefits of owning it, scarcity is a fair guide to its value; the warning applies to things you want to use. The post turns Cialdini's test into a sign of error.

The evidence is also less solid than the list suggests. Two items reach the post only through Cialdini's retelling: a conference paper, and an unpublished dissertation run by a businessman with the company's own sales staff. The post's hedged guess at the censorship source probably names a study of a different speech. The beef comparison group heard the usual sales pitch; nothing in Cialdini's account says they were told beef was plentiful. Other work puts limits on two of the findings. Censorship by an attractive expert did not raise interest in a speech the audience disagreed with. A 2006 meta-analysis reports that instructions to disregard evidence fail to remove its effect; its abstract does not say they generally increase it.

In short: an accurate retelling of Cialdini's examples that presents only as a fault what Cialdini calls a shortcut that usually works, and whose most striking example concerns a shortage that was real.
\end{afterword}
